\section{Introduction}
\label{sec:introduction}

Barge-in exposes a state-consistency problem that arises in interactive speech systems. When a user begins speaking while an assistant response is playing, language-model generation, text-to-speech (TTS) synthesis, and software playback may have advanced to different positions. The language model may already have generated a suffix that the synthesizer has only partly converted and the player has not consumed. Committing the complete generated response can make the next turn depend on information outside the observed playback path. Discarding the whole assistant turn instead loses both the retained prefix and its computation. Stopping audio is therefore only one part of barge-in handling: a cascaded system must also decide which assistant prefix remains committed and make every representation of the conversation agree with that decision.

This problem requires a precise account of progress. We distinguish \emph{software-consumed samples}, \emph{device-presented samples}, and \emph{acoustically heard content}. The first quantity is reported by playback software. The second additionally depends on application, operating-system, driver, and device buffers, while the third is further affected by transduction, propagation, and the listener's environment. This work observes only a software-consumed-sample cursor. Accordingly, ``playback-conditioned'' denotes state operations driven by this software cursor, not token-level knowledge of what reached the device or what a listener heard.

Cascaded automatic speech recognition (ASR)--large language model (LLM)--TTS systems remain attractive because their components can be selected, instrumented, and optimized separately. Recent work has addressed semantic triggering, incremental reasoning, interruption control, and streaming synthesis \citep{zou2026LTSVoiceAgent,mai2026RelayS2S,chen2025FireRedChat,du2024CosyVoice2}. Predictive ASR can also use partial recognition to prepare downstream results before the input is final \citep{schwarz2023PersonalizedPredictiveASR}. Within \emph{Speech Communication}, adjacent studies have examined low-latency integration for telephone-conversation processing \citep{morrone2024LowLatencyDiarization}, the effect of telecommunications latency on speaker transitions \citep{edwards2025TelecommunicationsLatency}, and conversation-history prompting for speech-LLM ASR \citep{zheng2026MultilevelContextualPrompting}. Faster or earlier computation, however, does not determine which generated content should survive a barge-in.

Playback-conditioned history truncation is established engineering practice. The OpenAI Realtime API accepts an \texttt{audio\_end\_ms} boundary and removes unplayed audio together with the associated transcript \citep{openaiRealtimeConversations}. Azure Voice Live provides \texttt{auto\_truncate}, which updates the previous response and session context following an interruption and documents a real-time-playback assumption in its estimate \citep{microsoftVoiceLiveAutoTruncation2026}. LiveKit exposes interruption semantics intended to align stored transcript history with spoken output \citep{livekitEventsErrorHandling}. We therefore do not claim the general principle that history should reflect played output. Nor do we treat cache shortening as novel, since Transformers already provides a \texttt{DynamicCache.crop} operation \citep{huggingFace2025TransformersCacheUtils,huggingFace2025KVCacheStrategies}.

The research object is narrower and cross-layer: an \emph{external-progress-conditioned joint prefix-state repair contract}. A software cursor is resolved through TTS-fragment metadata to a legal assistant commit boundary. That boundary is then applied consistently to the transformer key--value (KV) cache, attention mask, complete token ledger, assistant-content span, position progression, and chat role/end-of-turn (EOT) state. These components cannot be repaired independently. A cache of the intended length may still encode an invalid conversation if its mask, ledger, positions, or role boundary refers to a different prefix. Conversely, a transcript edit does not ensure that the retained model state represents that transcript.

The contract has four layers. First, \emph{boundary resolution} maps the software cursor through a TTS fragment to an assistant token span. Second, \emph{joint-state representation} treats the cache, mask, ledgers, content span, positions, and role/EOT state as one coupled object. Third, \emph{invariant-preserving transitions} define cropping and role reopening without duplicating or misclassifying structural EOT. Fourth, \emph{falsifiable validation} compares the production crop with direct layer-wise slicing of the same pre-crop snapshot, applies wrong-length negative controls, and checks matched recovery from identical retained states.

The article addresses five questions in decreasing order of evidential importance. RQ1 asks whether joint repair satisfies direct crop integrity and within-run matched-arm recovery exactness under a frozen model/backend, and what its fixed-protocol costs are. RQ2 compares software-cursor retention and complete-generation retention using fixed trajectories and automated information-reproduction probes. RQ3 characterizes how a soft trigger affects candidate availability, discarded-token ratio, and an optimistic synchronous-oracle latency bound. RQ4 audits full-string/full-prefill and segment-wise/incremental implementation paths while recognizing that they are not token-equivalent. RQ5 describes three interrupted-history naturalization implementations, whose present trajectories do not permit a causal comparison.

The sole core mechanism contribution is the implemented path
\begin{equation}
\begin{split}
\text{software cursor}&\rightarrow\text{TTS fragment}\rightarrow
\text{assistant token span}\\
&\rightarrow\text{KV crop}\rightarrow\text{role/EOT recovery}.
\end{split}
\label{eq:contribution-path}
\end{equation}
Under a frozen Qwen2-7B-Instruct, Transformers, BF16, and scaled dot-product attention (SDPA) configuration \citep{yang2024Qwen2}, the accepted exact-only evaluation covered 24 cases, 27 crop events, and 60 recovery steps. The retained production state was bitwise exact across 28 K/V layers against direct layer-wise slicing of the same snapshot, all wrong-length controls were detected, and matched recovery arms remained stepwise exact within the accepted run. These results support direct crop integrity and within-run matched-arm recovery exactness only. Earlier clean-reprefill protocols remain rejected; the accepted protocol does not establish clean-reprefill or continuation equivalence, or correctness across models, backends, or hardware.

The remaining studies provide supporting or exploratory evidence. A fixed-trajectory downstream probe produced uncertain differences across four automated operationalizations and is not a human-semantic or human--computer interaction evaluation. Candidate generation was evaluated in a synchronous segmented-text harness whose timestamps do not establish readiness before a real end of speech or production delivery. A separate implementation-path comparison was not token-equivalent. The paper consequently contributes a testable state contract and bounded software evidence rather than a claim of production barge-in latency or improved user experience.
